\section{Conclusion}
\label{sec:conclusion}

Penelitian ini mengevaluasi kerangka kerja fusi representasi laten Tri-Domain ViT (Face $\oplus$ Emotion $\oplus$ Age: 2,304 dimensi) untuk klasifikasi ras dan gender interseksional pada dataset DemogPairs. Hasil eksperimen menunjukkan bahwa fusi tri-domain menjadi konfigurasi terbaik pada 3 dari 4 classifier yang dievaluasi, mencakup SVM, LR, dan GNB. Dalam search space yang dievaluasi, model SVM teroptimasi dengan hyperparameter $C=10$ dan kernel polinomial derajat dua menghasilkan performa klasifikasi tertinggi, yakni Akurasi 93.70\% serta F1-Score 93.69\% (0.9369). Meskipun demikian, manfaat integrasi representasi multi-domain terbukti bersifat classifier-dependent karena RF justru meraih performa puncak pada skema dual-domain Emotion $\oplus$ Face dengan akurasi 86.85\%. Di samping itu, evaluasi kinerja subkelompok mengonfirmasi bahwa nilai F1-Score pada keenam subkelompok demografis terdistribusi stabil pada rentang 91.74\% hingga 96.14\% (0.9174 hingga 0.9614).

Terlepas dari capaian tersebut, penelitian ini memiliki sejumlah keterbatasan yang membuka peluang riset lanjutan. Pertama, cakupan demografis penyelidikan ini terbatas pada tiga kelompok ras makro (Asian, Black, dan White) pada dataset DemogPairs, sehingga temuan empiris yang diperoleh tidak dapat digeneralisasikan secara langsung ke populasi multirasial maupun kelompok etnisitas lain tanpa pengujian komprehensif tambahan. Kedua, ekstraksi representasi laten dijalankan secara terpisah (offline) memanfaatkan backbone ViT pra-latih yang dibekukan tanpa fine-tuning end-to-end secara simultan. Ketiga, evaluasi eksperimental masih berfokus pada citra lingkungan laboratorium terkontrol, sehingga respons model terhadap variasi pose ekstrem, oklusi, dan pencahayaan liar (in-the-wild) memerlukan pengujian lebih jauh. Oleh karena itu, arah penelitian mendatang disarankan untuk mengeksplorasi mekanisme fusi atensi adaptif lintas-modal (cross-attention transformer fusion) yang dilatih secara end-to-end untuk mempelajari bobot interaksi dinamis antardomain, serta memperluas pengujian pada dataset berskala besar yang lebih heterogen seperti FairFace atau UTKFace. Selain itu, penerapan metode visual explainability seperti Attention Rollout atau Grad-CAM sangat diperlukan untuk memetakan kontribusi spasial tiap domain secara transparan, disertai penegakan tata kelola etika kecerdasan buatan (Responsible AI) dengan pengawasan manusia (human-in-the-loop) untuk mencegah eksploitasi pada sistem pengawasan publik.
